\documentclass[12pt,twocolumn]{article}

\usepackage[T1]{fontenc}
\usepackage[utf8]{inputenc}
\usepackage{mathptmx}
\usepackage[margin=1.4cm]{geometry}
\usepackage{booktabs}
\usepackage{graphicx}
\usepackage{microtype}

\setlength{\columnsep}{0.6cm}
\setlength{\parskip}{0pt}
\setlength{\abovecaptionskip}{2pt}
\setlength{\belowcaptionskip}{0pt}
\setlength{\textfloatsep}{5pt}
\setlength{\intextsep}{5pt}
\renewcommand{\arraystretch}{0.9}
\pagestyle{empty}

\makeatletter
\renewcommand\section{\@startsection{section}{1}{\z@}
  {4pt plus 1pt minus 1pt}{2pt}{\normalfont\large\bfseries}}
\makeatother

\IfFileExists{tablas/cifras.tex}{\newcommand{\cifraRenglones}{147.411}
\newcommand{\cifraFacturas}{68.807}
\newcommand{\cifraProductos}{1.265}
\newcommand{\cifraCanasta}{2,14}
\newcommand{\cifraUnitarias}{28,203 (41.0\%)}
\newcommand{\cifraTicketMediano}{\$5,189}
\newcommand{\cifraTicketPromedio}{\$14,468}
\newcommand{\cifraChiCuadrado}{1313,2}
\newcommand{\cifraPValor}{1{,}5\times 10^{-280}}
\newcommand{\cifraFacturasPanel}{46,554}
\newcommand{\cifraEquivalentes}{8}
\newcommand{\cifraCombinaciones}{8}
\newcommand{\cifraSinCompletar}{12}
\newcommand{\cifraAprioriSinMemoria}{36}
\newcommand{\cifraAprioriCorridas}{60}
}{}
\providecommand{\cifraRenglones}{--}\providecommand{\cifraFacturas}{--}
\providecommand{\cifraProductos}{--}\providecommand{\cifraCanasta}{--}
\providecommand{\cifraUnitarias}{--}\providecommand{\cifraTicketMediano}{--}
\providecommand{\cifraTicketPromedio}{--}\providecommand{\cifraChiCuadrado}{--}
\providecommand{\cifraPValor}{--}\providecommand{\cifraFacturasPanel}{--}
\providecommand{\cifraEquivalentes}{--}\providecommand{\cifraCombinaciones}{--}
\providecommand{\cifraSinCompletar}{--}\providecommand{\cifraAprioriSinMemoria}{--}
\providecommand{\cifraAprioriCorridas}{--}

\newcommand{\incluirtabla}[1]{
  \IfFileExists{tablas/#1}{\begin{center}\footnotesize\input{tablas/#1}\end{center}}
  {\begin{center}\small\textit{(tabla pendiente: \texttt{make all})}\end{center}}}
\newcommand{\incluirfigura}[2]{
  \IfFileExists{../outputs/figs/#1}{
    \begin{figure}[h]\centering
    \includegraphics[width=0.96\linewidth]{../outputs/figs/#1}
    \caption{\footnotesize #2}
    \end{figure}}{
    \begin{center}\small\textit{(figura pendiente: \texttt{make all})}\end{center}}}

\title{\vspace{-1.3cm}\textbf{Reglas de asociación sobre transacciones de retail:
apriori frente a fp-growth}}
\author{Sara Albarracín}
\date{\vspace{-0.35cm}\small Laboratorio 03 --- 16 de octubre de 2026\vspace{-0.45cm}}

\begin{document}
\maketitle

\section*{1. Introducción}

Se analizan las facturas de una tienda de conveniencia con dos objetivos: describir el
comportamiento de compra y comparar dos implementaciones que resuelven el mismo
problema, apriori y fp-growth. El conjunto son 13 archivos DBF de Visual FoxPro, uno
por día, con \cifraRenglones{} renglones que se agrupan en \cifraFacturas{} facturas
sobre \cifraProductos{} productos. El dato que condiciona todo es el tamaño de la
canasta: \cifraCanasta{} renglones por factura y \cifraUnitarias{} de facturas con un
solo producto. Una canasta tan corta obliga a soportes mínimos muy bajos, y es ahí
donde las dos implementaciones dejan de costar lo mismo.

\section*{2. Datos y método}

Los archivos son DBF con una línea por \emph{line-item}, así que la factura se
reconstruye agrupando por archivo, caja, tipo de documento y consecutivo. La descripción
del producto es un \texttt{Varchar} de FoxPro, cuyo último byte guarda la longitud, y el
lector la interpreta al leerla. El procesamiento es un ETL de tres capas con 18
contratos de datos.
Los nombres se canonizan por código, porque 23 códigos tienen más de una descripción y
de otro modo el mismo producto repartiría su soporte entre dos ítems.

La fecha sale del nombre del archivo y la hora se estima repartiendo las facturas de cada
caja, en su orden real de facturación, sobre una jornada de 07:00 a 22:00. Se incluyen
los cinco tipos de documento; los que no son ventas son el 2\,\% de los documentos y el
55\,\% del valor, por lo que se cita ticket \emph{mediano} (\cifraTicketMediano{}) y no
promedio (\cifraTicketPromedio{}).

\section*{3. Resultados}

La demanda está muy concentrada: las empanadas aparecen en una de cada diez facturas,
el segundo producto del ranking ya baja al 41\,\% de ese volumen y el décimo no llega al
22\,\%. El ranking completo está en \texttt{outputs/tables/06\_top10.csv}.

Los 13 archivos cubren dos semanas desiguales; solo el tramo del 9 al 15 de marzo trae
los siete días una vez cada uno, y es el único donde comparar días es limpio. Sobre ese
panel de \cifraFacturasPanel{} facturas, el índice es el valor del día sobre el promedio
semanal.

\incluirtabla{dia_semana.tex}

El viernes concentra un 27\,\% más de tráfico que el promedio y el sábado un 21\,\%
menos ($\chi^2 = \cifraChiCuadrado$, $p = \cifraPValor$). La canasta, en cambio, es plana
en los siete días: cambia cuánta gente entra, no qué compra. Ese perfil ---creciente de
lunes a viernes, bajo el fin de semana--- corresponde a una tienda de entorno laboral o
escolar, no residencial.

Ambos algoritmos reciben la misma matriz binaria y los mismos parámetros, así que la
única variable es la implementación. De cada corrida se guarda una firma SHA-256 de sus
itemsets frecuentes: coinciden en \cifraEquivalentes{} de las \cifraCombinaciones{}
combinaciones comparables. Las otras \cifraSinCompletar{} no se pueden comparar porque
\textbf{apriori no terminó}: agotó la memoria en \cifraAprioriSinMemoria{} de sus
\cifraAprioriCorridas{} corridas, todas las de soporte igual o menor a 0,001.

\incluirtabla{benchmark.tex}

\incluirfigura{10_benchmark.png}{Tiempo y número de reglas según el soporte mínimo
(escalas logarítmicas). Mismo resultado, distinto costo.}

\section*{4. Respuestas al taller}

\textbf{1. Reglas predecibles ($k=1$).} Las de mayor confianza son de insumo, no de
consumo: caja de regalo $\rightarrow$ empaque (0,95). Entre productos de consumo, la
predecible es empanada $\rightarrow$ jugo.

\textbf{2. ¿Difieren los algoritmos?} No en el resultado: donde ambos terminan, los
itemsets son idénticos. Sí en costo, y de dos maneras. En tiempo, el orden se invierte
según el soporte: a 0,005 apriori es casi tres veces \emph{más rápido}, porque armar el
FP-tree tiene un costo fijo que no se amortiza cuando hay poco que minar; a 0,002
fp-growth ya es seis veces más rápido. En memoria la diferencia es cualitativa: apriori
materializa los candidatos de cada nivel, y con unos 500 ítems frecuentes eso son
125.000 pares contra 68.807 transacciones, de modo que el proceso muere. fp-growth no
genera candidatos y completa el mismo caso en 11 segundos.

\textbf{3. Efecto de los parámetros.} El soporte mínimo gobierna todo:

\incluirtabla{parametros.tex}

Bajar de 0,0005 a 0,0002 multiplica las reglas por 134, pero solo una de cada veinte
sobrevive al filtro de soporte absoluto. El soporte también decide qué algoritmo es
viable: por debajo de 0,001, apriori deja de ser una opción en esta máquina.
\texttt{max\_len} acota el crecimiento limitando el tamaño de los itemsets, aunque no
salva a apriori, porque el cuello está en los candidatos de nivel 2.

\textbf{4. Reglas espurias.} Sí, y con parámetros precisos: al soporte del análisis no
aparece ninguna, porque el umbral ya garantiza 34 facturas por regla. Hay que bajar a
0,0002 (14 facturas) para obtenerlas, y surgen reglas de lift en los miles sostenidas
por un puñado de facturas ---globos y serpentina $\rightarrow$ toper de cumpleaños,
lift 3.440 sobre 16 facturas---. El criterio aplicado contra ellas es el soporte
\emph{absoluto} ($\geq$30 facturas), que no depende del tamaño del conjunto.

\textbf{5. Reglas con $k>1$.} Existen, dominadas por combinaciones de armado de regalo.
Hay que deduplicar por itemset: cinco productos generan más de veinte reglas con
idénticos soporte y lift, que son la misma regla escrita de otra forma.

\textbf{6. Probabilidad al azar.} Si fueran independientes sería $P(A)\cdot P(B)$; el
cociente entre el soporte observado y ese valor es exactamente el \emph{lift}. Se
contrastó además la tabla $2\times2$ con chi-cuadrado y Fisher, para afirmar la
asociación con un $p$-valor.

\textbf{7. Reglas interesantes.} Las que unen soporte suficiente con lift sobre el azar
son de comida rápida más bebida: empanada $\rightarrow$ jugo (411 facturas), papas
$\rightarrow$ jugo, yogur $\rightarrow$ ponqué. Dos hallazgos fuera del patrón: la
tienda presta servicio de papelería (fotocopia $\rightarrow$ impresiones, lift 65), y
aparecen \textbf{sustitutos}, que el lift revela igual de bien ---sándwich
$\rightarrow$ empanada tiene lift 0,90---: un lift bajo 1 significa que los productos
compiten por la misma compra.

\textbf{8. Comportamiento del consumidor.} Todo apunta a consumo inmediato: canasta de
dos productos, ticket mediano de \cifraTicketMediano{}, 41\,\% de facturas de un solo
ítem y reglas que unen un alimento con una bebida. El cliente no hace mercado, resuelve
un hambre puntual, y que las reglas se mantengan estables entre días hábiles y fin de
semana confirma que cambia el volumen de visitas, no la lógica de la compra.

\section*{5. Conclusiones}

Elegir entre apriori y fp-growth no es cuestión de resultado sino de escala, y la
frontera no es el tiempo sino la memoria: producen los mismos itemsets, pero por debajo
de 0,001 apriori no llega a producirlos. A soporte alto, en cambio, apriori gana, así
que la recomendación no es universal: depende del régimen, y este conjunto ---canastas
de dos productos--- obliga al régimen donde fp-growth es la única opción. Ahí los
soportes altos no devuelven nada y los muy bajos sobre todo ruido; el soporte absoluto
resultó mejor filtro que el relativo. En lo descriptivo, la demanda se concentra en
pocos productos de consumo inmediato y el día mueve el tráfico pero no la conducta, lo
que sugiere ajustar inventario por día sin cambiar el surtido.

\end{document}
